\chapter[Lesson 9: Existence, Location and the Past]{Lesson 9: \ja{あります}/\ja{います}, Location Words and the Past Tense}

\section{Vocabulary}
\begin{center}
\begin{tabular}{lll@{\qquad}lll}
\toprule
\ja{つくえ} & tsukue & desk & \ja{いぬ} & inu & dog \\
\ja{いす} & isu & chair & \ja{ねこ} & neko & cat \\
\ja{テーブル} & tēburu & table & \ja{こども} & kodomo & child \\
\ja{へや} & heya & room & \ja{おとこのひと} & otoko no hito & man \\
\ja{ベッド} & beddo & bed & \ja{おんなのひと} & onna no hito & woman \\
\ja{えき} & eki & station & \ja{バスてい} & basutei & bus stop \\
\ja{こうえん} & kōen & park & \ja{びょういん} & byōin & hospital \\
\ja{スーパー} & sūpā & supermarket & \ja{ホテル} & hoteru & hotel \\
\bottomrule
\end{tabular}
\end{center}

\section{Grammar 1: \ja{あります} and \ja{います} ``there is''}

Japanese has two verbs of existence, chosen by what exists:
\begin{itemize}
  \item \ja{います} (\ja{いる}): living, self-moving things --- people, animals.
  \item \ja{あります} (\ja{ある}): inanimate things, plants, places, and \emph{events}.
\end{itemize}

\begin{center}
\begin{tabular}{llll}
\toprule
 & Affirmative & Negative & Example subject \\
\midrule
animate & \ja{います} & \ja{いません} & \ja{ねこ}, \ja{ともだち}, \ja{せんせい} \\
inanimate & \ja{あります} & \ja{ありません} & \ja{ほん}, \ja{き}, \ja{テスト} \\
\bottomrule
\end{tabular}
\end{center}

\subsection{Pattern A: place \ja{に} thing \ja{が} \ja{あります}/\ja{います}}
Introduces something new: ``In [place] there is [thing].''
\begin{example}
\leavevmode
\begin{itemize}
  \item \ja{へやにベッドがあります。} \emph{Heya ni beddo ga arimasu.}
        There is a bed in the room.
  \item \ja{こうえんにいぬがいます。} \emph{Kōen ni inu ga imasu.}
        There is a dog in the park.
  \item \ja{あしたテストがあります。} \emph{Ashita tesuto ga arimasu.}
        There is a test tomorrow. (event: \ja{あります})
\end{itemize}
\end{example}

\subsection{Pattern B: thing \ja{は} place \ja{に} \ja{あります}/\ja{います}}
Talks about where a known thing is: ``[Thing] is in [place].''
\begin{example}
\ja{わたしのかばんはつくえのうえにあります。}
\emph{Watashi no kaban wa tsukue no ue ni arimasu.} My bag is on the desk.\\
\ja{せんせいはどこにいますか。} \emph{Sensei wa doko ni imasu ka.} Where is the teacher?
\end{example}

\subsection{\ja{が}: a first look}
\ja{が} marks the subject when it is new information, and is the particle that goes with
\ja{あります}/\ja{います} in Pattern A. Question words as subjects always take \ja{が},
never \ja{は}: \ja{だれがいますか。} ``Who is there?''

\subsection{Possession}
\ja{あります} also expresses having things:
\ja{わたしはじてんしゃがありません。} ``I don't have a bicycle.''
\ja{きょうだいがいますか。} ``Do you have siblings?''

\section{Grammar 2: location words}

Location words are nouns. To say ``on the desk'' we say ``the desk's top'':
\ja{つくえのうえ}. Pattern: \textbf{X \ja{の} [location] \ja{に}}.

\begin{center}
\begin{tikzpicture}[x=1cm,y=1cm]
  \draw[thick, fill=brown!20] (-1.2,-0.8) rectangle (1.2,0.8);
  \node[font=\small] at (0,0) {X};
  \node[font=\small, align=center] at (0,1.6) {\ja{うえ}\\ue (on, above)};
  \node[font=\small, align=center] at (0,-1.6) {\ja{した}\\shita (under)};
  \node[font=\small, align=center] at (-2.8,0) {\ja{ひだり}\\hidari (left)};
  \node[font=\small, align=center] at (2.8,0) {\ja{みぎ}\\migi (right)};
  \node[font=\scriptsize, align=center] at (0,0.45) {\ja{なか} (inside)};
  \draw[-Stealth] (0,0.85) -- (0,1.15);
  \draw[-Stealth] (0,-0.85) -- (0,-1.15);
  \draw[-Stealth] (-1.25,0) -- (-1.95,0);
  \draw[-Stealth] (1.25,0) -- (1.95,0);
  % front/back depth
  \draw[dashed, gray] (1.2,0.8) -- (2.2,1.8);
  \node[font=\small, align=center] at (3.3,2.1) {\ja{うしろ} ushiro (behind)};
  \draw[dashed, gray] (-1.2,-0.8) -- (-2.2,-1.8);
  \node[font=\small, align=center] at (-3.1,-2.1) {\ja{まえ} mae (in front)};
  \node[font=\small, align=center] at (5.6,-1.2) {\ja{となり} tonari\\(next to, same kind)\\ \ja{ちかく} chikaku (near)\\ \ja{あいだ} aida (between)};
\end{tikzpicture}
\end{center}

\begin{center}
\begin{tabular}{lll}
\toprule
Word & Romaji & Meaning \\
\midrule
\ja{うえ} & ue & on, above \\
\ja{した} & shita & under, below \\
\ja{なか} & naka & inside \\
\ja{まえ} & mae & in front of \\
\ja{うしろ} & ushiro & behind \\
\ja{みぎ} & migi & right \\
\ja{ひだり} & hidari & left \\
\ja{となり} & tonari & next to (similar things side by side) \\
\ja{ちかく} & chikaku & near \\
\ja{AとBのあいだ} & A to B no aida & between A and B \\
\bottomrule
\end{tabular}
\end{center}

\begin{example}
\leavevmode
\begin{itemize}
  \item \ja{ねこはいすのしたにいます。} The cat is under the chair.
  \item \ja{ぎんこうはゆうびんきょくのとなりです。} The bank is next to the post office.
        (\ja{です} can replace \ja{にあります}.)
  \item \ja{えきのまえにバスていがあります。} There's a bus stop in front of the station.
  \item \ja{スーパーはびょういんとホテルのあいだにあります。} The supermarket is between
        the hospital and the hotel.
\end{itemize}
\end{example}

\section{Grammar 3: past tense of \ja{です}}

\begin{center}
\begin{tabular}{lll}
\toprule
 & Present & Past \\
\midrule
affirmative & \ja{がくせいです} & \ja{がくせいでした} (was a student) \\
negative & \ja{がくせいじゃないです} & \ja{がくせいじゃなかったです} (was not) \\
negative (formal) & \ja{がくせいではありません} & \ja{がくせいではありませんでした} \\
\bottomrule
\end{tabular}
\end{center}
\begin{example}
\ja{きのうはにちようびでした。} \emph{Kinō wa nichiyōbi deshita.} Yesterday was Sunday.\\
\ja{こどものとき、わたしはいいこどもじゃなかったです。} As a child I wasn't a good kid.
\end{example}

\section{Grammar 4: past tense of verbs}

Replace \ja{ます} with \ja{ました}; replace \ja{ません} with \ja{ませんでした}.
\begin{center}
\begin{tabular}{lllll}
\toprule
Verb & Present aff. & Present neg. & Past aff. & Past neg. \\
\midrule
\ja{たべる} & \ja{たべます} & \ja{たべません} & \ja{たべました} & \ja{たべませんでした} \\
\ja{いく} & \ja{いきます} & \ja{いきません} & \ja{いきました} & \ja{いきませんでした} \\
\ja{する} & \ja{します} & \ja{しません} & \ja{しました} & \ja{しませんでした} \\
\ja{ある} & \ja{あります} & \ja{ありません} & \ja{ありました} & \ja{ありませんでした} \\
\bottomrule
\end{tabular}
\end{center}

\begin{example}
\leavevmode
\begin{itemize}
  \item \ja{きのうとしょかんでべんきょうしました。} I studied at the library yesterday.
  \item \ja{せんしゅうのしゅうまつ、どこにもいきませんでした。} Last weekend I didn't go
        anywhere.
  \item \ja{パーティーにだれがいましたか。} Who was at the party?
\end{itemize}
\end{example}

\subsection{Past time words}
\ja{きのう} yesterday, \ja{おととい} the day before yesterday, \ja{せんしゅう} last week,
\ja{せんげつ} last month, \ja{きょねん} last year, \ja{〜まえ} \ldots ago
(\ja{みっかまえ} three days ago). Like \ja{きょう}, these take no \ja{に}.

\subsection{Duration: \ja{〜じかん}}
\ja{いちじかん} one hour, \ja{にじかん} two hours. No particle:
\ja{まいにちにじかんにほんごをべんきょうします。} ``I study Japanese two hours every day.''

\section{Kanji}
\begin{center}
\begin{tabular}{lllll}
\toprule
Kanji & Meaning & Readings & Example & \\
\midrule
\ja{人} & person & \ja{ジン}, \ja{ニン}; \ja{ひと} & \ja{日本人} & nihonjin, Japanese person \\
\ja{本} & book; origin & \ja{ホン}; \ja{もと} & \ja{本} & hon, book \\
\ja{上} & up, above & \ja{ジョウ}; \ja{うえ} & \ja{つくえの上} & tsukue no ue, on the desk \\
\ja{下} & down, below & \ja{カ}; \ja{した} & \ja{いすの下} & isu no shita, under the chair \\
\ja{中} & middle, inside & \ja{チュウ}; \ja{なか} & \ja{中国} & Chūgoku, China \\
\ja{山} & mountain & \ja{サン}; \ja{やま} & \ja{富士山} & Fujisan, Mt.\ Fuji \\
\ja{川} & river & \ja{セン}; \ja{かわ} & \ja{山川さん} & Yamakawa-san \\
\bottomrule
\end{tabular}
\end{center}

\begin{aside}
\ja{上}, \ja{下}, \ja{山}, \ja{川} are \emph{pictographs}: \ja{山} is three peaks, \ja{川}
is flowing water, and \ja{上}/\ja{下} point above/below a base line. \ja{日本} literally
means ``origin of the sun'', i.e.\ ``land of the rising sun''.
\end{aside}

\section*{Lesson 9 review}
\begin{itemize}
  \item \ja{います} for people/animals, \ja{あります} for things and events.
  \item Place \ja{に} thing \ja{が} \ja{あります}; thing \ja{は} place \ja{に} \ja{あります}.
  \item Location: X \ja{の} \ja{うえ}/\ja{した}/\ja{なか}/\ja{まえ}/\ja{うしろ}/\ja{となり}\ldots
  \item \ja{でした}, \ja{じゃなかったです}; verbs \ja{ました}, \ja{ませんでした}.
  \item Kanji: \ja{人本上下中山川}.
\end{itemize}
